% AUTORIDAD DE FUENTE:
% expose_VIII_body.tex, líneas 1976--2003; índices de escaneo 298--300;
% folios de fuente 287--289. Las correcciones materiales de la capa
% francesa se conservan reversiblemente mediante \sisi.

\medskip
\noindent\underline{6. Extensiones y biextensiones por el modelo de
Néron de} $\underline G_m$.

\medskip
\noindent 6.1. Sean $S$ un esquema, $T$ un subesquema cerrado y $Z$
un conjunto, que da un esquema constante truncado $Z_{S,T}$ sobre $S$
(5.7). Existe un homomorfismo canónico, funtorial en $Z$, de esquemas
étale sobre $S$,
\begin{equation}\tag{6.1.1}
 Z_S\longrightarrow Z_{S,T},
\end{equation}
que es un morfismo \underline{sobreyectivo}, como se desprende
inmediatamente de (5.7.1). Su funtorialidad en $Z$ muestra que, cuando
$Z$ es un grupo, (6.1.1) es un homomorfismo de esquemas en grupos
étale sobre $S$. Su núcleo es, por tanto, un esquema en grupos étale
sobre $S$, denotado por $Z_{S,U}$, y se tiene una sucesión exacta
canónica, funtorial en $Z$,
\begin{equation}\tag{6.1.2}
 0\longrightarrow Z_{S,U}\longrightarrow Z_S\longrightarrow Z_{S,T}
 \longrightarrow0.
\end{equation}
De (5.7.1) se deduce inmediatamente que $Z_{S,U}$ es la suma de una
familia de $S$-esquemas $S_n$, $n\in Z$, donde $S_n=S$ para $n=0$ y
$S_n=U$ para $n\ne0$. Esta descripción de $Z_{S,U}$ muestra fácilmente
que, para todo

%%%%%%%%%%  scan idx 299  |  folio 288  %%%%%%%%%%

grupo $F$ sobre $S$, la restricción a $U$ induce una biyección
\begin{equation}\tag{6.1.3}
 \mathrm{Hom}_{\mathrm{gr}}(Z_{S,U},F)
 \xrightarrow{\ \sim\ }
 \mathrm{Hom}_{\mathrm{gr}}(Z_U,F|U).
\end{equation}

\medskip
\noindent 6.1.4. En lo sucesivo suponemos que $Z$ es conmutativo, y
queremos examinar las extensiones de grupos conmutativos de $Z_{S,T}$
por un grupo conmutativo dado $G$ sobre $S$. Por la sucesión exacta
(6.1.2), la categoría de estas extensiones equivale a la categoría de
extensiones $F$ de $Z_S$ por $G$, provistas de una trivialización de la
imagen inversa de la extensión por $Z_{S,U}\to Z_S$. Por (6.1.3), tal
trivialización equivale a una trivialización de la extensión $F|U$ de
$Z_U$ por $G|U$.

En lo que sigue tomaremos como $Z$ el grupo $\mathbb Z$ de los
enteros. Así, $\mathbb Z_S$ es el grupo libre sobre un generador, y
VII 1.4 da una equivalencia entre la categoría de extensiones de
$\sisi{\mathbb Z_Z}{\mathbb Z_S}$ por $G$ y la categoría de torsors
bajo $G$. Aplicando el mismo resultado a $\mathbb Z_U$ y $G|U$ se
obtiene lo siguiente.

\medskip
\noindent\underline{Lema 6.2.} \underline{Sean $S$ un esquema, $T$ un
subesquema cerrado, $U=S\setminus T$, y $G$ un grupo conmutativo sobre
$S$. La categoría de grupos conmutativos sobre $S$ que son extensiones
de $\mathbb Z_{S,T}$} (definido en (5.7)) \underline{por $G$ equivale
a la categoría de pares $(F,s)$ formados por un torsor $F$ bajo $G$ y
una sección $s$ de $F|U$. Si $\alpha:S\to\mathbb Z_{S,T}$ es la
sección de $\mathbb Z_{S,T}$ imagen de la sección $1$ de $\mathbb Z_S$
sobre $\sisi{{}_{\overline G}S}{S}$, entonces el torsor correspondiente
a una extensión $\overline G$ de $\mathbb Z_{S,T}$ por $G$ es
$F=\alpha^*(\overline G)$ bajo $G$, provisto de la sección de $F|U$
dada por la restricción a $U$ de la sección unidad de $\overline G$.}

%%%%%%%%%%  scan idx 300  |  folio 289  %%%%%%%%%%
